\section{Emergence 的证据边界与工程含义}

到这里，课件暂停增加案例，转而总结 emergence 能支持什么、不能支持什么。它能说明继续扩展模型可能打开小模型实验无法预测的任务区域，也能提醒安全研究关注潜在 undesirable capabilities；但它尚不能给出精确阈值预测、无限规模极限或最佳架构配比。

\lecturefigure{slide-19.jpg}{Emergence 总结：尺度、数据、fine-tuning 与开放问题}{官方课件第 19 页；Wei et al., 2022}

\begin{knowledgebox}{读图：把总结拆成观察、推论与开放问题}
观察是：多项 few-shot abilities 和 prompting techniques 在被测试的模型家族中呈阈值行为。工程推论是：不要只用小模型判定方法上限，同时应探索数据与后训练把边界左移。开放问题则包括如何预测新能力、如何选择正确横轴、如何在大型预训练成本下做可信消融，以及 undesirable abilities 是否也会以类似方式出现。
\end{knowledgebox}

\teachervoice{讲者把 emergence 称为继续 scaling 的一个 implicit argument，因为昂贵扩展可能带来事先难以枚举的能力；但他也明确说预测未来涌现仍然很难，自己没有看到足够通用的方法。这不是“规模万能论”，而是关于未知收益与未知风险的研究理由。}

\subsection{大型预训练为什么难以形成干净理论}

学生追问应该怎样在 depth、attention heads、embedding width 之间分配新增参数。讲者回答，当时没有 principled recipe，很多选择受 TPU 容量和工程经验约束；更关键的是，完整 pretraining ablation 极其昂贵。因而我们常看到少量 model family 的 curve，却缺少实验设计中理想的单变量控制。

\begin{warningbox}{“没有做消融”不是“小问题”}
如果 data、tokenizer、architecture 与 optimizer 同时变化，某个阈值移动只能被解释为整个 recipe 的结果。把它简化成“参数更多导致能力出现”会过度归因。讲义、论文评审和产品声明都应把可观测关系与未经验证的机制分开。
\end{warningbox}

\subsection{Benchmark 也会老化}

BIG-Bench 的某些任务最初对 GPT-3/LaMDA 保持 flat，PaLM 出现后又有几十项超过 baseline。课堂甚至现场发现一个列表中的任务可能已被 ChatGPT 解决。Benchmark 因此不是静态难度尺：模型改进、prompting 方法和数据污染都可能让“hard”迅速失效。

\begin{importantbox}{一个可复用的 emergence 审计表}
\begin{enumerate}
  \item 是否给出 chance、human 或 prior-system baseline？
  \item 是否覆盖足够密的 model scales，并报告误差条或多 seed？
  \item 模型族之间是否同时改变 data、architecture 与 post-training？
  \item metric 是否连续，还是把连续质量阈值化成 exact match？
  \item 结果能否在新的 prompt、task split 和独立模型族上复现？
\end{enumerate}
\end{importantbox}

\subsection{本章小结}

Emergence 是值得测量的 empirical pattern，却不是完整因果理论。它最有价值的用途是暴露外推失败：小模型上无效的方法可能在大模型上有效，小模型上看似稳定的负趋势可能反转，旧 benchmark 也会被快速穿透。

